\section*{补充材料：数据、实现与完整结果}

\section{任务、数据与统计口径}
十四项任务标识为grasp、t1\_belt、t1\_occ、t1\_so3、t1\_xyz、t2\_se3、t2\_trans、t3\_full、t3\_phase、t4\_can、t4\_slider、t4\_sphere、t4\_sphere\_tabletop和t5\_replan。全部任务进入等权评分，包括零成功任务。正式可分发版本还需附实际评测器配置与数值事件阈值，不沿用旧稿中的历史抓取阈值代替。

共同训练池为11,200条；本方内部表征实验使用350条验证episode、每轮160,416个chunks。DreamVLA和OpenWAM保留各自1,400条heldout设置，因此350不能写成所有方法统一的验证规模。future screen的十五臂共18,900个有效结果，缺场与基础设施错误排除均为零。fixed140与补集分别报告；BASE2按完整1,400场选择配置。

令$b$为候选成功、参照失败的场景数，$c$为候选失败、参照成功的场景数，则差值为$100(b-c)/1260$个百分点。原分析对结果不一致的配对场景使用精确双侧McNemar检验，14项相对P00与13项相对P03共27项整体比较组成校正族，调整值为$\min(1,27p)$。378项逐任务检验只作探索性描述。源分析未提供置信区间，因此本文报告点估计和已有配对检验；这些检验不包含训练随机性。

\section{实现与预算}
D主干为四层、width768、8heads、factorized attention与split\_ffn，冻结DINOv3-L/16。当前图像为agentview224，196个当前视觉tokens。H0V没有视觉历史query，但保留20\,Hz的T8 EEF7及八个object tokens。对象统计包含位置、旋转向量、区间平移与旋转、时间间隔、线/角速度、线/角加速度及三类有效标志，来自GT。当前位姿条件屏蔽历史物体位姿及全部派生区间运动学。

实际位姿损失为有效标签上的平移平方误差加0.1倍SO(3)测地距离平方，外层权重单独设置。Action在训练归一化空间使用单位高斯源；scene flow在米制空间使用标准差0.02\,m的高斯源。位姿为确定性读出，不消费FM噪声和时间。旧稿的异方差NLL、对称性最小化、pose--flow一致性和不确定性bridge不属于这里已完成的模型。

P04的Wan目标由九帧RGB产生三个时间latent，不称为八个独立未来latent。P05为稠密二维累计位移，P06为真实codec块运动，P07为稀疏场景轨迹及可见性，P08为目标池化DINO，P09为目标轨迹。P08/P09是机制改编，不是整篇方法复现。Aalpha为P03增加权重0.5的运动模长MSE；P12、P13、P14分别加depth、DINO、动态RGB，P15按对应单项权重合加，不能识别全部两两交互。

ACT和DP各50,000更新、global32，从头训练；$\pi_{0.5}$从pi05\_base开始，30,000更新、global256；DynamicVLA为50,000更新、global44。OpenWAM从官方Alpha基座重新微调五轮，共34,490更新、global192；DynamicWAM用官方full初始化，40,000更新、global128。DreamVLA独立20轮与原40轮配方分别48,220与96,440更新，global512。执行器也不同：ACT用temporal aggregation，DP执行8步，$\pi_{0.5}$执行10步，DynamicWAM执行16步，本方内部screen执行1步。不同系统的结果不能解释为等预算主干优越性。

\section{Q1完整比较的边界}
已完成当前图像、Window、Codec-MV、Flow和PACT-Age；Flow12D、Wan、PACT-Uniform和fresh GT历史参照尚未形成完整表。PACT-Age将56个历史query按远至近的$[2,4,6,8,10,12,14]$分配；Uniform为七组各8个，均为逐帧query压缩改编。已完成的五臂各有350 episode、1,050窗口的探针提取与306项有效读出；这些数组还需整理为可报告的探针效应。

HISTORY1采用冻结DINO-L、原Alpha，学习率第一轮10,026更新cosine后，三轮保持$2\times10^{-5}$；BASE2使用DINO-S/CNN、Alpha关闭和覆盖40,104更新的cosine。需要在共同配方下完成当前位姿/GT运动与当前帧/视觉历史对照。匹配当前状态的历史替换、遮挡持续时间与采样间隔诊断，需同时保留当前状态匹配证据及实际时间。\pending{完整九臂、匹配输入核查、探针统计与闭环干预。}

\section{全部未来监督与组合结果}
下列英文标识表与独立英文补充材料采用相同数据，直接由研究仓库结果摘要生成；它们是论文导出，不是新的实验事实源。场景流包含目标区域，不能将P01当作纯context-only对照。正向点估计、不利结果和任务失败全部保留。

% Generated from Research source summaries; do not edit numerical cells.
\begin{table*}[!htp]
\centering\small
\begingroup
\setlength{\tabcolsep}{3pt}
\renewcommand{\arraystretch}{1.10}
\caption{All fifteen future-supervision conditions, grouped by the objective changed. Each uses one training seed, four epochs, common privileged object-motion input, and action-only deployment. The 140 development scenes and 1,260-scene complement remain separate. Scene flow includes target regions.}
\label{tab:all_future}
\begin{tabular*}{\textwidth}{@{\extracolsep{\fill}}llrrr@{}}
\toprule
 & & Development & \multicolumn{2}{c}{Disjoint complement} \\
\cmidrule(lr){3-3}\cmidrule(lr){4-5}
ID & Future supervision & Success / 140 & Success / 1,260 & SR (\%) \\
\midrule
\addlinespace[2pt]
\multicolumn{5}{@{}l}{\emph{Reference and component objectives}} \\[-1pt]
P00 & No future supervision & 88 & 791 & 62.78 \\
P01 & 3-D scene flow & 95 & 782 & 62.06 \\
P02 & Target pose & 92 & 789 & 62.62 \\
P03 & Scene flow + target pose & 93 & 786 & 62.38 \\
\addlinespace[2pt]
\multicolumn{5}{@{}l}{\emph{Scene replacements: target-pose objective retained}} \\[-1pt]
P04 & Wan latent + pose & 92 & 783 & 62.14 \\
P05 & Dense 2-D displacement + pose & 87 & 746 & 59.21 \\
P06 & Codec motion + pose & 86 & 795 & 63.10 \\
P07 & Sparse scene tracks + pose & 90 & 807 & 64.05 \\
\addlinespace[2pt]
\multicolumn{5}{@{}l}{\emph{Target replacements: scene-flow objective retained}} \\[-1pt]
P08 & Scene flow + object DINO & 89 & 794 & 63.02 \\
P09 & Scene flow + target tracks & 88 & 808 & 64.13 \\
\addlinespace[2pt]
\multicolumn{5}{@{}l}{\emph{Additional objectives: scene flow + target pose retained}} \\[-1pt]
$A_{\alpha}$ & P03 + motion magnitude & 89 & 778 & 61.75 \\
P12 & P03 + depth & 90 & 791 & 62.78 \\
P13 & P03 + DINO & 92 & 788 & 62.54 \\
P14 & P03 + dynamic RGB & 89 & 790 & 62.70 \\
P15 & P03 + depth / DINO / dynamic RGB & 94 & 802 & 63.65 \\
\bottomrule
\end{tabular*}
\endgroup
\end{table*}

% Generated from Research source summaries; do not edit numerical cells.
\begin{table*}[!htp]
\centering\small
\begingroup
\setlength{\tabcolsep}{3pt}
\renewcommand{\arraystretch}{1.10}
\caption{Paired outcomes over the same 1,260 scenes. Gain means candidate success/reference failure; loss means candidate failure/reference success. Intervals are pointwise task-stratified paired bootstrap intervals for the SR difference, conditional on the trained checkpoints (Appendix~A). Exact two-sided McNemar tests use the full family of fourteen P00 and thirteen P03 contrasts: $p_{\mathrm{adj}}=\min(1,27p)$. P03 against itself is inapplicable.}
\label{tab:paired_future}
\begin{tabular*}{\textwidth}{@{\extracolsep{\fill}}lrrrrrrrr@{}}
\toprule
 & \multicolumn{4}{c}{P00: no future supervision} & \multicolumn{4}{c}{P03: scene flow + target pose} \\
\cmidrule(lr){2-5}\cmidrule(lr){6-9}
ID & Gain & Loss & 95\% CI (pp) & Adjusted $p$ & Gain & Loss & 95\% CI (pp) & Adjusted $p$ \\
\midrule
\addlinespace[2pt]
\multicolumn{9}{@{}l}{\emph{Reference and component objectives}} \\[-1pt]
P01 & 80 & 89 & $[-2.70, 1.27]$ & 1 & 72 & 76 & $[-2.14, 1.51]$ & 1 \\
P02 & 61 & 63 & $[-1.83, 1.51]$ & 1 & 61 & 58 & $[-1.43, 1.90]$ & 1 \\
P03 & 64 & 69 & $[-2.14, 1.35]$ & 1 & --- & --- & --- & --- \\
\addlinespace[2pt]
\multicolumn{9}{@{}l}{\emph{Scene replacements: target-pose objective retained}} \\[-1pt]
P04 & 74 & 82 & $[-2.54, 1.27]$ & 1 & 71 & 74 & $[-2.06, 1.59]$ & 1 \\
P05 & 53 & 98 & $[-5.40, -1.75]$ & 0.00844 & 47 & 87 & $[-4.92, -1.43]$ & 0.0188 \\
P06 & 81 & 77 & $[-1.59, 2.22]$ & 1 & 68 & 59 & $[-1.03, 2.46]$ & 1 \\
P07 & 85 & 69 & $[-0.63, 3.17]$ & 1 & 81 & 60 & $[-0.16, 3.49]$ & 1 \\
\addlinespace[2pt]
\multicolumn{9}{@{}l}{\emph{Target replacements: scene-flow objective retained}} \\[-1pt]
P08 & 75 & 72 & $[-1.59, 2.06]$ & 1 & 59 & 51 & $[-0.95, 2.22]$ & 1 \\
P09 & 91 & 74 & $[-0.56, 3.25]$ & 1 & 79 & 57 & $[0.00, 3.49]$ & 1 \\
\addlinespace[2pt]
\multicolumn{9}{@{}l}{\emph{Additional objectives: scene flow + target pose retained}} \\[-1pt]
$A_{\alpha}$ & 60 & 73 & $[-2.78, 0.71]$ & 1 & 53 & 61 & $[-2.30, 1.03]$ & 1 \\
P12 & 73 & 73 & $[-1.83, 1.83]$ & 1 & 63 & 58 & $[-1.27, 2.06]$ & 1 \\
P13 & 70 & 73 & $[-2.06, 1.59]$ & 1 & 60 & 58 & $[-1.51, 1.83]$ & 1 \\
P14 & 71 & 72 & $[-1.90, 1.75]$ & 1 & 62 & 58 & $[-1.35, 2.06]$ & 1 \\
P15 & 74 & 63 & $[-0.95, 2.70]$ & 1 & 63 & 47 & $[-0.32, 2.94]$ & 1 \\
\bottomrule
\end{tabular*}
\endgroup
\end{table*}

% Generated from Research source summaries; do not edit numerical cells.
\begin{table*}[!htp]
\centering\small
\begingroup
\setlength{\tabcolsep}{2.1pt}
\renewcommand{\arraystretch}{1.10}
\caption{Future-study successes: ninety scenes per task and 1,260 per condition. Columns follow the task families in Table~\ref{tab:tasks}; all tasks are retained, including zero-success tasks. Occl. denotes physical occlusion; Trans. denotes translating-receptacle placement; Tabletop denotes tabletop sphere capture.}
\label{tab:future_tasks}
\begin{tabular*}{\textwidth}{@{\extracolsep{\fill}}lrrrrrrrrrrrrrrr@{}}
\toprule
 & Reference & \multicolumn{4}{c}{T1: target motion} & \multicolumn{2}{c}{T2: placement} & \multicolumn{2}{c}{T3: passage} & \multicolumn{4}{c}{T4: capture} & T5 & \\
\cmidrule(lr){2-2}\cmidrule(lr){3-6}\cmidrule(lr){7-8}\cmidrule(lr){9-10}\cmidrule(lr){11-14}\cmidrule(lr){15-15}
ID & Grasp & Belt & Occl. & SO(3) & XYZ & SE(3) & Trans. & Full & Phase & Can & Slider & Sphere & Tabletop & Replan & Total \\
\midrule
\addlinespace[2pt]
\multicolumn{16}{@{}l}{\emph{Reference and component objectives}} \\[-1pt]
P00 & 85 & 90 & 90 & 75 & 76 & 0 & 8 & 3 & 15 & 90 & 69 & 45 & 86 & 59 & 791 \\
P01 & 84 & 90 & 90 & 64 & 69 & 0 & 18 & 0 & 11 & 90 & 74 & 46 & 90 & 56 & 782 \\
P02 & 87 & 90 & 90 & 78 & 73 & 0 & 7 & 0 & 0 & 90 & 71 & 59 & 89 & 55 & 789 \\
P03 & 86 & 90 & 88 & 75 & 77 & 0 & 7 & 0 & 6 & 86 & 80 & 50 & 87 & 54 & 786 \\
\addlinespace[2pt]
\multicolumn{16}{@{}l}{\emph{Scene replacements: target-pose objective retained}} \\[-1pt]
P04 & 83 & 85 & 84 & 73 & 76 & 0 & 20 & 0 & 0 & 90 & 72 & 54 & 89 & 57 & 783 \\
P05 & 84 & 89 & 90 & 66 & 75 & 0 & 10 & 0 & 0 & 90 & 71 & 23 & 88 & 60 & 746 \\
P06 & 86 & 88 & 89 & 79 & 70 & 0 & 6 & 4 & 2 & 89 & 72 & 60 & 85 & 65 & 795 \\
P07 & 82 & 90 & 90 & 78 & 71 & 0 & 17 & 0 & 0 & 89 & 79 & 55 & 90 & 66 & 807 \\
\addlinespace[2pt]
\multicolumn{16}{@{}l}{\emph{Target replacements: scene-flow objective retained}} \\[-1pt]
P08 & 84 & 90 & 90 & 71 & 81 & 0 & 13 & 0 & 0 & 90 & 77 & 60 & 90 & 48 & 794 \\
P09 & 84 & 90 & 89 & 65 & 78 & 0 & 21 & 0 & 0 & 90 & 77 & 59 & 90 & 65 & 808 \\
\addlinespace[2pt]
\multicolumn{16}{@{}l}{\emph{Additional objectives: scene flow + target pose retained}} \\[-1pt]
$A_{\alpha}$ & 86 & 90 & 90 & 72 & 75 & 0 & 3 & 0 & 0 & 85 & 73 & 55 & 90 & 59 & 778 \\
P12 & 83 & 90 & 90 & 76 & 76 & 0 & 5 & 0 & 0 & 89 & 77 & 57 & 90 & 58 & 791 \\
P13 & 86 & 90 & 90 & 80 & 73 & 0 & 6 & 0 & 0 & 88 & 78 & 54 & 90 & 53 & 788 \\
P14 & 85 & 89 & 90 & 72 & 79 & 0 & 5 & 0 & 0 & 90 & 80 & 53 & 88 & 59 & 790 \\
P15 & 87 & 90 & 90 & 79 & 81 & 0 & 11 & 0 & 0 & 90 & 75 & 54 & 90 & 55 & 802 \\
\bottomrule
\end{tabular*}
\endgroup
\end{table*}

% Generated from Research source summaries; do not edit numerical cells.
\begin{table*}[t]
\centering\small
\caption{Native evaluator failure categories out of 1,260 episodes per condition. Safety is the evaluator category, not a certified safety measure. All failures stay in the score denominator.}
\label{tab:future_failures}
\begin{tabular}{@{}lrrrrr@{}}
\toprule
ID & Boundary & Safety & Task violation & Timeout & All failures \\
\midrule
P00 & 84 & 297 & 17 & 71 & 469 \\
P01 & 62 & 324 & 23 & 69 & 478 \\
P02 & 77 & 331 & 17 & 46 & 471 \\
P03 & 76 & 295 & 51 & 52 & 474 \\
P04 & 79 & 310 & 41 & 47 & 477 \\
P05 & 87 & 356 & 14 & 57 & 514 \\
P06 & 96 & 250 & 23 & 96 & 465 \\
P07 & 47 & 329 & 20 & 57 & 453 \\
P08 & 61 & 322 & 42 & 41 & 466 \\
P09 & 67 & 299 & 35 & 51 & 452 \\
Aalpha & 57 & 315 & 65 & 45 & 482 \\
P12 & 62 & 327 & 40 & 40 & 469 \\
P13 & 71 & 326 & 27 & 48 & 472 \\
P14 & 70 & 322 & 18 & 60 & 470 \\
P15 & 59 & 343 & 19 & 37 & 458 \\
\bottomrule
\end{tabular}
\end{table*}

% Generated from Research source summaries; do not edit numerical cells.
\begin{table*}[t]
\centering\small
\begingroup
\setlength{\tabcolsep}{2.5pt}
\renewcommand{\arraystretch}{1.10}
\caption{Combination-study successes, grouped by future target and visual encoder. DINO is frozen DINO-S; CNN is trained from scratch. Minus/plus denotes current-target RGB reconstruction off/on; parentheses retain source condition IDs. Each task has one hundred scenes, including development scenes, and each condition has one training seed. Totals use the predefined 1,400-scene selection pool.}
\label{tab:combination_tasks}
\begin{tabular*}{\textwidth}{@{\extracolsep{\fill}}lrrrrrrrrrrrr@{}}
\toprule
 & \multicolumn{4}{c}{P15: additional targets} & \multicolumn{4}{c}{P07: scene tracks} & \multicolumn{4}{c}{P09: target tracks} \\
\cmidrule(lr){2-5}\cmidrule(lr){6-9}\cmidrule(lr){10-13}
Task & \makecell{DINO$-$\\(B)} & \makecell{DINO$+$\\(A)} & \makecell{CNN$-$\\(D)} & \makecell{CNN$+$\\(C)} & \makecell{DINO$-$\\(H)} & \makecell{DINO$+$\\(G)} & \makecell{CNN$-$\\(J)} & \makecell{CNN$+$\\(I)} & \makecell{DINO$-$\\(L)} & \makecell{DINO$+$\\(K)} & \makecell{CNN$-$\\(N)} & \makecell{CNN$+$\\(M)} \\
\midrule
\addlinespace[2pt]
\multicolumn{13}{@{}l}{\emph{Reference grasp}} \\[-1pt]
Grasp & 90 & 89 & 91 & 92 & 90 & 92 & 87 & 90 & 79 & 85 & 77 & 87 \\
\addlinespace[2pt]
\multicolumn{13}{@{}l}{\emph{T1: target motion}} \\[-1pt]
Belt & 100 & 100 & 99 & 100 & 100 & 100 & 100 & 97 & 99 & 100 & 90 & 99 \\
Occlusion & 100 & 100 & 99 & 99 & 99 & 100 & 100 & 97 & 96 & 99 & 82 & 98 \\
SO(3) & 86 & 71 & 80 & 88 & 75 & 87 & 69 & 86 & 69 & 62 & 64 & 78 \\
XYZ & 58 & 58 & 82 & 75 & 73 & 75 & 89 & 82 & 44 & 56 & 82 & 84 \\
\addlinespace[2pt]
\multicolumn{13}{@{}l}{\emph{T2: receptacle motion}} \\[-1pt]
SE(3) & 0 & 0 & 0 & 0 & 0 & 0 & 0 & 0 & 0 & 0 & 0 & 0 \\
Translation & 8 & 4 & 6 & 5 & 9 & 1 & 14 & 4 & 4 & 5 & 9 & 14 \\
\addlinespace[2pt]
\multicolumn{13}{@{}l}{\emph{T3: constrained passage}} \\[-1pt]
Full & 0 & 0 & 0 & 0 & 0 & 0 & 0 & 0 & 0 & 0 & 0 & 0 \\
Phase & 0 & 0 & 0 & 0 & 0 & 0 & 28 & 72 & 0 & 0 & 0 & 2 \\
\addlinespace[2pt]
\multicolumn{13}{@{}l}{\emph{T4: object capture}} \\[-1pt]
Can & 99 & 99 & 99 & 98 & 100 & 99 & 99 & 100 & 97 & 99 & 96 & 99 \\
Slider & 88 & 85 & 88 & 91 & 74 & 83 & 86 & 81 & 84 & 84 & 88 & 88 \\
Sphere & 71 & 67 & 67 & 63 & 47 & 40 & 53 & 56 & 68 & 68 & 62 & 64 \\
Tabletop sphere & 100 & 100 & 100 & 99 & 99 & 100 & 95 & 98 & 100 & 100 & 100 & 100 \\
\addlinespace[2pt]
\multicolumn{13}{@{}l}{\emph{T5: post-impact replanning}} \\[-1pt]
Replan & 73 & 63 & 64 & 57 & 74 & 72 & 76 & 72 & 80 & 72 & 79 & 57 \\
\midrule
Total / 1,400 & 873 & 836 & 875 & 867 & 840 & 849 & 896 & 935 & 820 & 830 & 829 & 870 \\
\bottomrule
\end{tabular*}
\endgroup
\end{table*}


P05相对P00和P03的校正$p$分别为0.00844和0.01877。P06的safety类失败最少，为250次，但timeout最多，为96次；P15的timeout最少，为37次，safety仍有343次。减少某一类失败不能单独说明策略更优。BASE2候选I相对J有817场共同成功、118场仅I成功、79场仅J成功、386场共同失败，净差39；该选择仍为单seed开发结果。

直接检验预测粒度，需要在同一兼容表示内匹配身份信息、物理时域及token/成本，再改变角色与精度分配。现有跨表示替换不能代替该对照。生成预测与接触事件应按同一物理时间对齐，不能用未来FM验证loss替代自由生成误差，也不能把预测误差下降直接写成动作机制。\pending{匹配粒度对照、预测—实际—接触图及未来oracle诊断。}

\section{D/E、七配置信息流与交叉确认}
BACKBONE0当前只有准备与CPU接线，无GPU训练、SR或成本结果。两臂计划各fresh seed0、1、2；若沿用四轮，共六次训练240,624更新、8,400场评测。活动主干参数差目标约5\%；相同曝光或参数量不等于相同FLOPs。既定选择规则使用2个百分点的实际重要差异和非劣界限，以及10\%的p95成本收益条件；只有区间支持对应判据时才称优越、等效或非劣。三seed不明确时，两臂按既有方案共同扩至五seed。

信息流研究的R0/R1区分辅助监督，R1两种部署区分是否额外生成未被读取的future，R1--R4区分访问关系，R2/R5与R3/R6区分同步和顺序程序。三seed对应21个训练配置、24个部署配置，按每端点1,400场计33,600场，均为计划数。Alpha固定为零。scene-flow生成状态形状为$[B,S,196,8,3]$；若最终选择其他未来表示，需要先完成生成接口。顺序模型的第二阶段GT教师强制与预测部署失配，以及5次joint对5+5次顺序调用的成本差异，都进入解释。

可替换未来接口建立后，再进行oracle、匹配swap、时间偏移或校准噪声诊断。动作变化只说明敏感性，有用预测还需由任务表现与安全结果共同验证。改变信息流后回访表示，并将所选信息流用于落选主干；必要多seed与未参与选择场景确认完成后，交付固定$(H,Z,B,G)$供预训练使用。\pending{上述配置与交叉比较的实测结果。}

\section{真实数据基座、规模与部署}
固定设计后研究既有loss权重、训练配方及架构内的深度/宽度；新增表示或信息路径仍属于设计变更。DROID准备检查时间、标定、深度、命令与实测运动、夹爪语义和训练集归一化。目标pose需额外估计并验证，仅生产最终所选目标与必要消融资产，保留无效标签和动作的原因与掩码。

正式基座由真实数据预训练，仿真与真机从同一基座分别后训练。对有无预训练，既比较共同下游预算，也比较双方等预算调参后的表现，单列预训练成本。数据规模、模型规模与计算量分别变化；未运行配置不进入scaling曲线。最终阶段复用已完成规模与迁移结果，仅补外部泛化、静态兼容性、真实传送带与实际延迟证据。旧稿180次实机是历史预算，不能当作完成数或当前已采用样本量。

部署计时应区分采集、可用、预处理、推理、排队、下发与应用，报告硬件、精度、warm-up、样本数和同步方式。当前公开策略调用范围不包含外部I/O与oracle，不能标为完整环境端到端时延。\pending{正式基座和scaling、独立后训练、原生外部成绩、真机次数和端到端计时。}

\section{难任务、夹爪与定性证据}
T3-full的800条示教释放只在$h_7$监督，$h_0$缺失。尾窗口短训有2/16成功，随后选定场景复核1/6，不能据此宣称模型已稳定掌握完整任务。T3-phase的三seed纠偏实验在256更新时为0/96、12/96、26/96，1,024更新时均为0/96；首动作加权六端点共576场全零，GT角色诊断另六端点576场全零。将Action初始源置零使三个256更新端点变为0/96、1/96、46/96，这是观察已有结果后选择的探索干预。

夹爪“最近三态”将指令投影到距离最近的$\{-1,0,+1\}$标签，边界为$\pm0.5$。第五轮端点的1,260场计数分别为sign698、固定0.1死区742、最近三态749、历史保持741。最近三态未显示明确总体优势，历史保持还增加其他任务失败，均未成为生产默认。第四至第八轮更广泛的死区分析尚未通过严格历史轨迹一致性审计，因此未列入主表改进。

正文两幅概念图明确不表示真实rollout，第三图只绘制实测成功率。待补定性图需对齐同一场景、相机、查询时刻、预测时刻、实际状态与接触事件，包含不利球体条件、保留轨迹候选及有效的未来到动作干预。视频需标明场景、时长、播放速度、剪辑、重置与人工干预；不以单个成功片段替代完整分母。\pending{同步预测/rollout资产与匿名视频。}
